\section{Proofs}\label{app:proofs}
% Translated from G4/PROOFS.md (self-reviewed; independently reviewed 26/10/2026: 0 critical, 4 major, 12 minor, all
% addressed). Keep in sync with that file. Citations marked VERIFY in refs.bib must be checked (G4-6) before submission.

\paragraph{Assumptions.} (A0) $\W\sim p(\W\mid X,y)$, independent of $(x_*,y_*)$ given $(X,y)$; nuisance parameters of
the observation model, if any, are a priori independent of $(\phi,X,y)$. (A1) $\tfm(\cdot\mid x,D)=\ppd(\cdot\mid x,D)$ for
every context the algorithm queries. (A2) $\ppd(k\mid x,D)>0$ for all $k$ and every context of positive probability.
(A3) $\lik$ is known and $\sum_y p(y\mid X)\lik(y)>0$. For instance-level observations
$\lik(y)=\prod_i\lik_i(y_i)$; the set of free points is $F=\{i:\lvert\{k:\lik_i(k)>0\}\rvert\ge2\}$.

\subsection{Lemma 1}
\begin{proof}[Proof of \cref{lem:fullcond}]
Given $\phi$ the pairs $(x_j,y_j)$ are i.i.d., hence exchangeable, so the joint law of $(D_{-i},(x_i,y_i))$ equals that of
$(D_{n-1},(x_*,y_*))$, a context of $n-1$ points and a query. The information $(X,y_{-i})$ coincides with $(x_i,D_{-i})$.
Therefore
\[p(y_i\mid y_{-i},X)=\frac{p((x_i,y_i),D_{-i})}{\sum_kp((x_i,k),D_{-i})}=p(y_i\mid x_i,D_{-i}).\qedhere\]
\end{proof}
The PPD conditions on the query features $x_*$ as well, which may carry information about $\phi$; this is the quantity
the PFN training loss targets, since $x_*$ is drawn from the prior together with the context \citep{muller2022pfn}.

\begin{corollary}[Positivity]\label{cor:positivity}
Under (A2), $p(y\mid X)>0$ for every $y\in[C]^n$.
\end{corollary}
\begin{proof}
Some $y^0$ has $p(y^0\mid X)>0$. Change one coordinate at a time, $y^0\to y^1\to\dots\to y^m=y$, step $k$ changing
point $s_k$. If $p(y^{k-1}\mid X)>0$, then $p(y^{k-1}_{-s_k}\mid X)>0$, the context $D^{k-1}_{-s_k}$ has positive
probability, and by \cref{lem:fullcond} and (A2)
$p(y^k\mid X)=p(y^{k-1}_{-s_k}\mid X)\,p(y^k_{s_k}\mid x_{s_k},D^{k-1}_{-s_k})>0$. Induct on $k$.
\end{proof}

\subsection{Theorem 1}
\begin{proof}[Proof of \cref{thm:exact}(a)]
The update of point $i$ is $P_i(y,y')=\pi_i(y'_i\mid y_{-i})\mathbb 1[y'_{-i}=y_{-i}]$ with
$\pi_i(k\mid y_{-i})\propto\tfm(k\mid x_i,D_{-i})\lik_i(k)$. By \cref{lem:fullcond} and (A1),
$\pi_i(k\mid y_{-i})\propto p(k\mid y_{-i},X)\lik_i(k)\propto\pi(k,y_{-i})$, the full conditional of $\pi$; hence
$\pi(y)P_i(y,y')=\pi(y_{-i})\pi_i(y_i\mid y_{-i})\pi_i(y'_i\mid y_{-i})$ is symmetric in $(y,y')$. By
\cref{cor:positivity}, $\pi(y)>0$ iff $\lik(y)>0$, a product set; any two of its points are joined by changing one
coordinate at a time inside the set, each step with probability at least $\pi_i(\cdot\mid\cdot)/\lvert F\rvert>0$, so the
random-scan kernel $P=\lvert F\rvert^{-1}\sum_{i\in F}P_i$ is irreducible, and aperiodic since $P(y,y)>0$. A finite,
irreducible, aperiodic chain converges geometrically. The implementation sweeps the free points in a uniformly random
order, $P_{\mathrm{sw}}=\frac1{\lvert F\rvert!}\sum_\sigma P_{\sigma(1)}\cdots P_{\sigma(\lvert F\rvert)}$; each factor
preserves $\pi$, and $P_{\mathrm{sw}}$ is reversible because the adjoint of a product is the product in reverse order and
the set of permutations is closed under reversal.
\end{proof}

\begin{proof}[Proof of \cref{thm:exact}(b)]
Choose a bag uniformly; if it is not pure, choose uniformly an ordered pair $(i,j)$ in it with $y_i\neq y_j$ and propose
$y'$ by swapping their labels. Brook's lemma through the intermediate state $(y'_i,y_{-i})$ gives
\[
\frac{p(y'\mid X)}{p(y\mid X)}=\frac{p(y'_i\mid y_{-i},X)}{p(y_i\mid y_{-i},X)}\cdot\frac{p(y'_j\mid y'_i,y_{-ij},X)}{p(y_j\mid y'_i,y_{-ij},X)},
\]
and \cref{lem:fullcond} with (A1) turns the right-hand side into the ratio $R_{ij}$ of four TFM predictives that the
algorithm computes with two forward passes. A swap preserves the counts, so $\lik(y')=\lik(y)$ and $R_{ij}=\pi(y')/\pi(y)$,
independently of the order of $(i,j)$. The number of differently labelled pairs depends only on the counts, so the
proposal is symmetric and the chain is Metropolis with stationary distribution $\pi$. Two assignments of a bag with
the same counts are joined by a sequence of swaps preserving the counts, each with positive probability by
\cref{cor:positivity}; bags are independent, so the product chain is irreducible (the lazy version is also aperiodic).
\end{proof}

\begin{proof}[Proof of \cref{thm:exact}(c)]
(Inductive.) The ergodic theorem for a finite irreducible chain applied to the bounded function
$y\mapsto p(y_*\mid x_*,(X,y))$, with (A1), gives almost sure convergence of \eqref{eq:predict} to $\bar p$.
By the law of total probability and (A0), $p(y_*\mid x_*,X,\W)=\sum_yp(y\mid x_*,X,\W)\,p(y_*\mid x_*,(X,y))$; the
difference of the expectations of one $[0,1]$-valued function under two measures is at most their TV distance.
If the prior is separable, $\phi=(\phi_x,\phi_{y\mid x})$ independent a priori with $p(x\mid\phi)=p(x\mid\phi_x)$ and
$p(y\mid x,\phi)=p(y\mid x,\phi_{y\mid x})$, then $p(y\mid x_*,X,\W)\propto\lik(y)\int p(y\mid X,\phi_{y\mid x})p(\phi_{y\mid x})d\phi_{y\mid x}$
times a factor that does not depend on $y$, so $p(y\mid x_*,X,\W)=p(y\mid X,\W)$ and the bound is zero.
(Transductive.) Adding $x_*$ as a free point with $\lik_*\equiv1$ gives a chain for $p(y,y_*\mid X,x_*,\W)$ that is exact by
part (a) (for LLP, the swap chain combined with single-site updates of $x_*$, each preserving the target). By
\cref{lem:fullcond}, the Rao--Blackwell term is $\tfm(y_*\mid x_*,(X,y^{(s)}))$, whose average converges to
$p(y_*\mid x_*,X,\W)$.
\end{proof}

\begin{proof}[Proof of \cref{prop:bayes}]
By the tower property, $\E[S(f,y_*)]=\E\big[\E[S(f,y_*)\mid X,\W,x_*]\big]$. For each value of $(X,\W,x_*)$, $f$ is a fixed
distribution and $y_*\sim p(\cdot\mid x_*,X,\W)$; strict propriety gives the unique minimiser
$f=p(\cdot\mid x_*,X,\W)$. Its risk is finite (at most $\log C$ for the log score).
\end{proof}

\subsection{Theorem 2}
\begin{proof}[Proof of \cref{thm:stability}]
\emph{Existence and uniqueness.} The state space is finite, so $\tilde P$ has a stationary distribution. With $\tfm>0$
every single-site update and every MH acceptance probability is positive, so $\tilde P$ has the transition graph of $P$
and is irreducible; $\tilde\pi$ is unique. Compatibility of the conditionals of $\tfm$ is never used.
\emph{Perturbation identity.} Let $\Delta=\tilde\pi-\pi$ and $\nu=\tilde\pi(\tilde P-P)$, so $\nu(\Omega)=0$. From
$\tilde\pi=\tilde\pi\tilde P$ and $\pi=\pi P$, $\Delta=\nu+\Delta P$, and iterating, $\Delta=\sum_{t<T}\nu P^t+\Delta P^T$.
\emph{Contraction.} For a signed measure $\mu$ with $\mu(\Omega)=0$, $\lVert\mu K\rVert_1\le\delta(K)\lVert\mu\rVert_1$.
If $\tau(P)<\infty$ then $\delta(P^T)\to0$, the remainder vanishes, and $\lVert\Delta\rVert_1\le\lVert\nu\rVert_1\tau(P)$; if
$\tau(P)=\infty$ there is nothing to prove (no aperiodicity is needed). Finally
$\lVert\nu\rVert_1\le\sum_y\tilde\pi(y)\lVert\tilde P(y,\cdot)-P(y,\cdot)\rVert_1=2\bar\varepsilon$.
\emph{Single-site error.} $\tilde P(y,\cdot)-P(y,\cdot)$ averages the differences of the single-site updates, and each is the
TV distance between $\tfm\lik_i$ and $\ppd\lik_i$ after normalisation. With $Z_p=\sum_k\ppd(k)\lik_i(k)$,
$2\,\TV(\tfm_\lik,\ppd_\lik)\le\big(\sum_k\lik_i(k)\lvert\tfm(k)-\ppd(k)\rvert+\lvert Z_q-Z_p\rvert\big)/Z_p\le4(\max_k\lik_i(k)/Z_p)\,\TV(\tfm,\ppd)$,
so the error is at most $2\kappa\varepsilon_{\TV}$ with $\kappa=\sup\max_k\lik_i(k)/Z_p$.
\emph{Prediction.} $\lvert\E_{\tilde\pi}\tfm-\E_\pi\ppd\rvert\le\lvert\E_{\tilde\pi}(\tfm-\ppd)\rvert+\lvert\E_{\tilde\pi}\ppd-\E_\pi\ppd\rvert\le\varepsilon'_{\TV}+\TV(\tilde\pi,\pi)$.
\end{proof}
\emph{Relation to classical bounds.} With $E=\tilde P-P$ and the fundamental matrix $Z$, $\tilde\pi-\pi=\nu Z$
\citep{schweitzer1968,meyer1980}, which gives $\TV(\tilde\pi,\pi)\le\tau_1(Z)\bar\varepsilon$ with $\tau_1(Z)\le\tau(P)$
\citep{seneta1988,cho2001}; \eqref{eq:stability} is the weaker, more interpretable form. It is attained by a two-state
chain. For a single-site sweep the per-sweep error is at most $\lvert F\rvert$ times the per-update error; the MH swap
kernel has error at most $4\varepsilon_\infty$, with $\varepsilon_\infty=\sup\lvert\log\tfm-\log\ppd\rvert$.

\subsection{Block updates}
\begin{proposition}[Blocks]\label{prop:blocks}
(a) A block update whose proposal is the product of block-wise TFM predictives (for LLP, conditioned on bag counts),
accepted with the MH ratio computed by Brook's lemma along the changed labels, leaves $\pi$ invariant when $\tfm=\ppd$.
(b) Without the MH step, \cref{thm:stability} applies with an error $\bar\varepsilon_B$ that is bounded only under an
extra stability assumption on the PPD, and the bound grows with $n$ for fixed $K$.
\end{proposition}
\begin{proof}[Proof of \cref{prop:blocks}]
(a) The proposal $Q_B$ depends only on $y_{-B}$, which the step does not change, so the update is MH-within-Gibbs with
an independent proposal. Brook's lemma along the changed labels with \cref{lem:fullcond} gives the exact ratio
$p(y'\mid X)/p(y\mid X)$, multiplied by the ratio of $\lik$ (one for LLP). The proposal is positive on the support of the
block and acceptance is positive, so the chain is irreducible.
(b) Let $P_{\mathrm{bG}}$ be exact block Gibbs. Under the stability assumption that adding one labelled point or one
unlabelled feature to a context changes the PPD by at most $\beta$ in TV, the chain rule and a hybrid argument give
$\TV\big(\prod_{i\in B\cap F}\ppd(\cdot\mid x_i,D_{-B}),\,p(y_{B\cap F}\mid y_{-(B\cap F)},X)\big)\le m_f(m-1)\beta$
for a block of $m$ points of which $m_f$ are free, and after reweighting by $\lik$,
$\bar\varepsilon_B\le2\kappa_B\,m_f(m-1)\beta$; \cref{thm:stability} with $P=P_{\mathrm{bG}}$ concludes.
\end{proof}
\emph{Why this is vacuous at our sizes.} With $\beta=O(1/n)$ and fixed $K$ ($m=n/K$), $m_f(m-1)\beta=O(n/K^2)$ grows
with $n$. For bag-blocks, $\kappa_B$ is the inverse probability that all bags of a group match their counts under the
product proposal: for $n\approx2{,}000$ and $K=10$, $\log_{10}\kappa_B\approx15.8$ for bags of 8. On an exactly solvable
LLP problem (3 bags of 4), the uncorrected bag-block chain has $\TV(\tilde\pi,\pi)$ between 0 and about 0.6 depending on
the prior and grouping, so the approximation is real at small $n$; at our sizes we measure its effect on predictions
instead (\cref{sec:ablation}).

\subsection{When the naive TFM suffices}
\paragraph{Complement to \cref{prop:naive}(a).} For binary labels with class prior $\varpi$, for any score $s$,
$\mathrm{AUC}(\eta)-\mathrm{AUC}(s)\le\E\lvert s-\tilde\eta\rvert/(\lambda\varpi(1-\varpi))$: the AUC regret of a learner fitted
to naive labels is its own estimation error, amplified by $1/\lambda$.
\begin{proof}[Proof of \cref{prop:naive} and of the complement]
(a) With $P(\tilde w=k\mid x,y)=M_{y,k}$ independent of $x$ and $M_{y,k}=b_k$ for $y\neq k$,
$\tilde\eta_k=M_{k,k}\eta_k+b_k(1-\eta_k)=b_k+\lambda_k\eta_k$, $\lambda_k=M_{k,k}-b_k$. For PLL with each wrong class
entering the candidate set independently with probability $q$ and the naive label chosen uniformly from the set,
$M_{y,y}=a=\E[1/(1+N)]=\frac{1-(1-q)^C}{Cq}$ with $N\sim\mathrm{Bin}(C-1,q)$ and $b=(1-a)/(C-1)$, so $\lambda=(Ca-1)/(C-1)>0$
for $q<1$. For LLP with random bags of size $k$, $P(\tilde w_i=1\mid x_i,y_i)=(y_i+(k-1)\varpi)/k$, so $\lambda=1/k$.
A strictly increasing affine map preserves ranks, hence one-vs-rest AUC. For the binary bound, write
$\varpi(1-\varpi)\mathrm{AUC}(s)=\E[\eta(X)(1-\eta(X'))h_s(X,X')]$ for independent copies, with
$h_s=\mathbb 1[s(X)>s(X')]+\frac12\mathbb 1[s(X)=s(X')]$; symmetrising,
$\varpi(1-\varpi)[\mathrm{AUC}(\eta)-\mathrm{AUC}(s)]=\frac12\E\big[\lvert\Delta\rvert(\mathbb 1[s\text{ reverses the pair}]+\frac12\mathbb 1[s\text{ ties},\Delta\neq0])\big]$
with $\Delta=\eta(X)-\eta(X')$. With $t=(s-b)/\lambda$, which has the ranks of $s$, a reversed or tied pair has
$\lvert\Delta\rvert\le\lvert\eta(X)-t(X)\rvert+\lvert\eta(X')-t(X')\rvert$, so the right-hand side is at most
$\E\lvert s-\tilde\eta\rvert/\lambda$. For $\lambda=1$ and scores that are class-probability estimates, this is Corollary~12 of \citet{agarwal2014}.
(b) $p_\theta(\tilde w)=\sum_Op_\theta(O)R(\tilde w\mid O)$ gives
$\partial_\theta\log p_\theta(\tilde w)=\E_\theta[\partial_\theta\log p_\theta(O)\mid\tilde w]$, and Jensen's inequality for
conditional expectations gives $I_{\tilde w}(\theta)\le I_O(\theta)$, with equality iff the score of $O$ is
$\sigma(\tilde w)$-measurable.
\end{proof}
\emph{Information lost by the naive labels} (logistic model $\eta=\sigma(\theta x)$, $\theta=2$, $x\sim\mathcal N(0,I)$; per
point). LLP: the Godambe information of the pointwise (composite) likelihood on naive labels, relative to the
information of the bag counts, is 0.134, 0.060, 0.033, 0.017, 0.009 for bags of 2, 4, 8, 16, 32. PLL: the naive labels
retain 0.24--0.26 of the information in the candidate sets at $q=0.3$ and 0.06--0.10 at $q=0.6$ ($C\in\{3,4,7\}$).
For PU-SCAR and label noise the ratio is 1.
